\chapter{Visão Geral}

\section{O jogo}

\textbf{Void Descent} é um \emph{roguelike} top-down em tempo real construído em
JavaScript puro com renderização WebGL. Cada partida regenera todo o conteúdo do jogo
(mapas, atlas de tiles, áudio, comportamento de inimigos) por algoritmos em tempo de
execução, sem nenhum \emph{asset} pré-fabricado.

A premissa é descer três andares procedurais, \textit{Antecâmara}, \textit{Subsistemas}
e \textit{Núcleo}, eliminando inimigos sala por sala e coletando armas e
\textit{upgrades}, até derrotar o boss final. A morte é permanente; cada \emph{run} é
uma nova execução do gerador.

\section{Equipe}

\begin{itemize}
  \item Luis Gustavo Olímpio
  \item Lauan Matheus
  \item José Melquíades
  \item João Leonardi
\end{itemize}

\section{Stack tecnológico}

\begin{itemize}
  \item \textbf{JavaScript ES6+}, sem \textit{frameworks}, ESLint ou bundlers.
    Carregamento direto de \texttt{<script>} em ordem.
  \item \textbf{WebGL 1.0} com \textit{shaders} de vértice e fragmento personalizados, e
    iluminação por vértice de até 32 luzes pontuais simultâneas.
  \item \textbf{Sonant-X}, \textit{port} JavaScript do sintetizador Sonant
    (Marcus Geelnard, Nicolas Vanhoren), gera todo o áudio em tempo de \emph{boot}.
  \item \textbf{Web Audio API}: \texttt{AudioContext} e \texttt{GainNode} mestre para
    \textit{routing} unificado e controle de \emph{mute}.
  \item \textbf{Electron 32} para empacotamento como \texttt{.exe} \textit{portable}
    Windows x64, via \texttt{electron-builder}.
\end{itemize}

\section{Métricas do projeto}

\begin{tabular}{lr}
  \toprule
  \textbf{Métrica} & \textbf{Valor} \\
  \midrule
  Linhas de código (JS+HTML) & $\sim$2.100 \\
  Tamanho do código fonte & 77 KB \\
  Quantidade de arquivos JS & 8 \\
  Tamanho do \texttt{.exe} \textit{portable} & 73 MB \\
  \quad Chromium embutido & $\sim$60 MB \\
  \quad Node runtime & $\sim$10 MB \\
  \quad Bundle do jogo (asar) & $\sim$80 KB \\
  Tempo médio de \emph{run} & 8 a 12 minutos \\
  \bottomrule
\end{tabular}

O contraste entre o código fonte (77 KB) e o \texttt{.exe} final (73 MB) decorre da
escolha de empacotamento via Electron, conforme especificação do GDD para garantir
experiência \emph{nativa} de \textit{desktop}. A geração procedural reside na camada do
\emph{bundle}, não no \emph{wrapper} de distribuição.

\section{Estrutura desta documentação}

A documentação está dividida em nove capítulos, cada um focando uma dimensão ortogonal
do projeto:

\begin{enumerate}
  \item \textbf{Visão Geral}, este capítulo.
  \item \textbf{Arquitetura de Componentes}, com a divisão modular e as dependências
    (\textit{component diagram}).
  \item \textbf{Sistema de Entidades}, com a hierarquia de classes do mundo do jogo
    (\textit{class diagram}).
  \item \textbf{Máquina de Estados}, com as transições do estado global
    (\textit{state diagram}).
  \item \textbf{Game Loop} \textit{micro} (sala-a-sala) e \textit{macro} (run completa),
    representado por \textit{activity diagrams}.
  \item \textbf{Sistemas Procedurais}: geração de \emph{dungeon}, síntese de áudio e
    \textit{tinting} de atlas (\textit{sequence diagrams}).
  \item \textbf{Build e Deploy}, com o \emph{pipeline} de empacotamento e distribuição
    (\textit{deployment diagram}).
  \item \textbf{Backlog Técnico, Decisões e Aprendizados}, com o caminho real do
    projeto: o que escolhemos, o que descartamos e o que aprendemos.
  \item \textbf{Considerações Finais}, com o balanço do que foi entregue, o que
    faríamos diferente e o que vem depois.
\end{enumerate}
